%! TEX root = main.tex

범위 데이터에 대한 안전한 검색 시스템은 데이터의 기밀성, 무결성, 그리고 검색 효율성 사이의 균형을 맞추기 위해 다양한 암호학적 기법과 시스템 구조를 통해 연구되어 왔다. 본 장에서는 기존 연구들을 주요 기술적 접근 방식에 따라 분류하고, 본 연구와의 차별성을 논한다.

\subsection{암호화 검색 및 순서 보존 암호}
초기 연구들은 암호화된 상태에서 데이터를 검색하기 위해 암호문의 순서를 평문의 순서와 동일하게 유지하는 순서 보존 암호(Order-Preserving Encryption, OPE)를 제안하였다. Agrawal et al. \cite{agrawal2004}은 수치 데이터에 대한 OPE 기법을 처음으로 제안하였으며, Shen et al. \cite{shen2021}은 이를 발전시켜 상태를 발생시키지 않는(Stateless) OPE 기법을 통해 아웃소싱된 데이터베이스에서의 범위 검색 효율성을 높였다. 또한, Popa et al.\cite{popa2011}이 제안한 CryptDB는 양파 암호화(Onion Encryption)라는 계층 구조를 통해 SQL 쿼리 유형에 따라 OPE나 결정론적 암호화(DET) 등 다양한 암호학 기술을 선택적으로 적용하여 실용적인 암호화 데이터베이스를 구현하였다. 

\bigskip

그러나 이러한 기법들은 암호문의 속성을 일부 노출한다는 본질적인 취약점이 존재한다. Naveed et al.\cite{naveed2015}은 OPE로 암호화된 데이터베이스에서 노출된 순서 및 빈도 정보를 바탕으로 원본 데이터가 복구될 수 있음을 증명하였다. 나아가 Islam et al.\cite{islam2012}은 일반적인 검색 가능 암호화 시스템에서도 노출되는 데이터 접근 패턴만으로 추론 공격을 통해 민감한 정보가 유출될 수 있음을 경고하였다. 본 연구는 이러한 보안 취약점들을 원천적으로 차단하기 위해 데이터의 순서나 접근 패턴을 포함해, 쿼리의 속성을 제외한 그 어떤 부가적인 정보도 노출하지 않는 영지식 증명(zk-SNARK)을 활용하며, 데이터 분포 노출을 막기 위해 사전 정의된 범위만을 증명하는 방식을 채택한다.

\subsection{공간 인덱스 및 비대칭 내적 연산 기반 검색}
앞선 순서 보존 암호의 데이터 순서 노출 문제를 방어하기 위해, 데이터의 평문 대소 관계를 숨기면서도 거리 비교를 가능하게 하는 특수 암호 연산과 공간 인덱스를 결합한 연구들이 등장하였다. 대표적으로 Wang et al.\cite{wang2013}은 다차원 범위 검색을 위해 R-tree의 최소 경계 사각형(MBR)을 비대칭 스칼라 내적 보존 암호(ASPE)로 암호화한 $\hat{R}$-tree를 제안하였다. 이 시스템은 쿼리와 데이터를 비대칭적으로 암호화하여, 서버가 원문을 모르는 상태에서도 내적 연산의 부호만으로 검색 범위 포함 여부를 판별할 수 있게 한다. 그러나 이 방식은 본 연구가 타겟하는 의료 데이터 검색 시스템에는 근본적으로 적용할 수 없는 치명적인 구조적 한계를 지닌다.

\bigskip

첫째, $\hat{R}$-tree의 MBR은 사용자가 임의로 지정하는 것이 아니라, 입력된 실제 데이터들의 최소값과 최대값을 둘러싸는 형태로 동적으로 결정된다. 반면 수치형 의료 데이터는 세계보건기구(WHO) 가이드라인과 같이 변하지 않는 '절대적인 임상 참고치'를 기준으로 진단과 검색이 이루어져야 하므로, 데이터의 분포에 따라 기준선이 매번 달라지는 동적 공간 인덱스 방식은 본 연구의 요구사항을 애초에 충족시킬 수 없다.

\bigskip

둘째, 순서 정보 은닉을 위해 데이터를 암호화된 노드 단위로 버킷화(Bucketization)하기 때문에, 쿼리 시 실제 조건에 맞지 않는 데이터까지 함께 반환되는 오탐(False Positives)이 필연적으로 발생한다. 의료 데이터 도메인에서는 미세한 수치 차이가 완전히 다른 임상적 진단을 의미하므로, 오탐을 허용하는 검색 모델은 의료 데이터 유통에 있어 가장 중요한 '결과의 신뢰성'을 상실하게 만든다. 이 밖에도 쿼리 처리를 위해 서버가 암호화된 인덱스를 순회하며 지속적으로 내적 연산을 수행해야 하는 연산 부담 역시 존재한다. 결론적으로, 동적인 범위 설정과 오탐을 수용하는 기존의 공간 인덱스 기반 검색은 정밀성이 요구되는 수치형 의료 데이터 도메인에 적용하기 부적합하다.

\bigskip

\subsection{동형 암호 및 학습형 인덱스 기반 검색}
공간 인덱스의 한계를 넘어, 데이터의 지역성을 기계학습으로 모델링하고 이에 동형 암호를 결합하여 프라이버시를 강화한 연구도 제안되었다. Wang et al.\cite{wang2025}이 제안한 SLRQ 시스템은 기계학습 기반의 학습형 인덱스를 구축하고, 이를 Paillier 가산 동형 암호로 암호화하여 아웃소싱하는 방식을 채택하였다. 이 시스템은 학습된 데이터 분포를 기반으로 후보군을 빠르게 탐색한 뒤, 오탐을 제거하기 위해 동형 암호화된 상태에서 개별 튜플과 쿼리 값을 정밀 비교하는 연산을 수행한다. 그러나 이와 같은 동형 암호 기반의 검색 방식은 실시간 시스템에 적용하기에 다음과 같은 몇가지 한계를 지닌다.

\bigskip

첫째, 동형 암호의 기하급수적인 연산 오버헤드로 인한 실시간성 훼손이다. 암호화된 상태에서 정밀한 대소 비교를 수행하기 위해서는 막대한 암호문 프로토콜(SIC, SM 등) 연산이 동반된다. 시스템의 보안성을 높이기 위해 암호화 키 길이를 증가시킬 경우 이러한 연산 지연은 기하급수적으로 증가하며, 결과적으로 단일 쿼리 처리에 수 초 이상의 시간이 소요되어 대규모 트래픽을 처리하는 실시간 검색 시스템에 부적합하다.

\bigskip

둘째, 복호화 권한의 존재로 인한 잠재적 정보 유출 위험성이다. SLRQ는 무거운 동형 암호 연산의 부하를 줄이기 위해 보조 연산 서버에 데이터베이스의 전체 비밀키를 제공하고, 난수가 섞인 형태이긴 하나, 데이터와 쿼리를 중간에 평문으로 복호화하여 비교하는 과정을 거친다. 이는 근본적으로 제3의 서버가 원문을 복원할 수 있는 권한과 수단을 갖게 됨을 의미하며, 해당 서버가 침해당하거나 노이즈 생성 패턴이 분석될 경우 의료 데이터의 기밀성이 치명적으로 유출될 수 있는 위험을 내포한다.

\bigskip

셋째, 동적 데이터 업데이트에 대한 취약성이다. 학습형 인덱스는 고정된 데이터 분포에 의존하기 때문에, 지속적으로 새로운 임상 데이터가 누적되는 의료 환경에서는 데이터 분포의 왜곡이 발생한다. 이를 해결하기 위해서는 정기적으로 전체 인덱스 모델을 재학습하고 다시 암호화하여 업로드해야 하는 막대한 유지보수 비용이 발생한다.

\subsection{검증 가능한 데이터베이스}
데이터의 기밀성 뿐만 아니라 쿼리 결과의 무결성을 보장하기 위한 연구들도 활발히 진행되었다. Zhang et al.\cite{zhang2015}의 IntegriDB와 vSQL은 아웃소싱된 데이터베이스에서 SQL 쿼리 결과가 정확하고 완전함(Sound \& Complete)을 암호학적으로 증명하는 시스템을 제안하였다. Li et al.\cite{li2023}의 ZKSQL은 이를 영지식 증명으로 확장하여 데이터 프라이버시를 강화하였다.

\bigskip

그러나 이러한 시스템들은 쿼리 시점에 실시간으로 증명을 생성해야 하므로 연산 오버헤드가 매우 크다. 특히 ZKSQL의 경우 일반적인 SQL 쿼리에 대해 수 초 내지는 수 분 이상의 증명 생성 시간이 소요되어, 실시간성이 요구되는 검색 서비스에 적용하기에는 한계가 있다. 본 연구는 이러한 실시간 증명 생성의 병목을 해소하기 위해, 데이터 업로드 시점에 '사전 정의된 범위'에 대한 증명을 미리 생성해 두고 검색 시에 이를 단순 조회하는 방식을 제안하여 검색 성능을 O(1) 수준(DB 조회 시간)으로 획기적으로 단축하였다.

\subsection{하드웨어 및 블록체인 기반 접근}
소프트웨어적 접근의 성능 한계를 극복하기 위해 신뢰 실행 환경(TEE)이나 블록체인을 활용하는 시도도 있었다. Zhou et al. \cite{zhou2021}의 VeriDB는 Intel SGX와 같은 신뢰 하드웨어 내에서 쿼리를 실행하여 무결성을 보장하며, Xu et al.\cite{xu2019}의 vChain은 블록체인 데이터베이스 상에서 누산기(Accumulator) 기반의 ADS(Authenticated Data Structure)를 통해 불리언 범위 쿼리의 무결성을 검증한다. 

\bigskip

하지만 하드웨어 기반 접근은 특정 제조사의 CPU에 의존해야 하고, 특정 연산이 특정 하드웨어에서 처리될 때 발생하는 여러 부산물(캐시 접근 패턴, 연산 소요 시간, 전력 소비량 , 전자기파 등)을 분석하여 진행되는 사이드 채널 공격(Side-channel Attack)의 위험이 있다. vChain과 같은 블록체인 기반 시스템은 데이터의 기밀성 보장보다는 무결성 검증에 초점이 맞춰져 있어 민감한 의료 데이터 보호에는 부족함이 있다. 본 연구는 특정 하드웨어에 의존하지 않고 범용적인 zk-SNARK 알고리즘만으로 기밀성과 무결성을 동시에 달성한다.

\bigskip

이처럼, 기존의 연구들은 검색 효율성을 위해 보안을 일부 희생하거나(OPE, CryptDB), 보안을 위해 효율성을 희생하는(ZKSQL, IntegriDB) 트레이드오프를 지니고 있다. 본 연구는 "zk-SNARK 기반의 사전 증명 생성"이라는 새로운 접근법을 통해, 데이터 원문과 순서 정보를 완벽히 은닉하면서도 실용적인 검색 속도를 보장하는 새로운 절충안을 제시한다. 그리고 동시에 현실의 데이터 수요를 고려해, 다중 쿼리의 처리에 대한 데이터 프라이버시와 소유권 증명을 모두 만족하는 방법론을 제시한다.
